% Data-derived report section; editorial comments are preserved in this template.
\section{Simulation Results}
Table~\ref{tab:performance} compares the circuits at 200~MHz IF. The passive
conversion loss is the negative of the signed gain entry. The active mixer
trades DC power for conversion gain and lower loaded noise figure. The
passive mixer avoids a DC transconductance bias, but its finite switch
conductance and limited LO swing reduce signal transfer. Its small
compression advantage does not imply a higher third-order intercept;
compression and intermodulation characterize different nonlinear behavior.

\begin{table}[!t]
\caption{Performance at 200 MHz IF}
\label{tab:performance}
\centering\footnotesize
\setlength{\tabcolsep}{4pt}
\begin{tabularx}{\columnwidth}{@{}Xrr@{}}
\toprule
Parameter & Active & Passive \\
\midrule
Voltage gain $G_V$ (dB) & 8.57 & -5.98 \\
Transducer gain $G_T$ (dB) & 1.58 & -12.96 \\
Loaded SSB NF (dB) & 11.15 & 16.69 \\
Loaded DSB NF (dB) & 8.12 & 13.71 \\
IIP3 (dBm/tone) & -2.72 & -5.47 \\
Input P1dB (dBm) & $-15.55\pm0.08$ & $-15.08\pm0.08$ \\
Core DC power (mW) & 0.685 & 0 \\
\bottomrule
\end{tabularx}
\end{table}

\subsection{Waveforms and Spectra}
The active output extrema are 0.622 and 1.010~V; capacitive LO feedthrough
can briefly push an output above the DC supply. The shared source node of
the LO pair spans 0.196--0.305~V. These settled extrema are more informative
for headroom than the static bias solution alone. The passive outputs span
approximately $-22.9$ to 22.2~mV. The differential output spectra show the
200~MHz IF together with out-of-band feedthrough.


\reportfigure{active_transient.png}{Active mixer: settled RF input and differential output, including feedthrough.}{fig:active-transient}

\reportfigure{active_spectrum.png}{Active mixer: differential output spectrum at $-55$~dBm input; full-band and IF views, with 50~MHz DFT bins.}{fig:active-spectrum}

\reportfigure{passive_transient.png}{Passive mixer: settled RF input and differential output, including feedthrough.}{fig:passive-transient}

\reportfigure{passive_spectrum.png}{Passive mixer: differential output spectrum at $-55$~dBm input; full-band and IF views, with 50~MHz DFT bins.}{fig:passive-spectrum}

\subsection{Compression, Intermodulation and Convergence}
The compression sweeps and two-tone plots use available source power per
tone. The horizontal compression reference is the mean of the first three
low-power gains minus 1~dB. The accepted cubic interval resolves both IM3
products above numerical residuals. Table~\ref{tab:checks} summarizes slope
and convergence checks; the dominant numerical sensitivity in this campaign
is the active mixer's tighter-repeat IIP3 change.

\begin{table}[!t]
\caption{Extraction and convergence checks}
\label{tab:checks}
\centering\footnotesize
\setlength{\tabcolsep}{4pt}
\begin{tabularx}{\columnwidth}{@{}Xrr@{}}
\toprule
Parameter & Active & Passive \\
\midrule
Fundamental slope (mean) & 0.996 & 0.992 \\
IM3 slope (mean) & 3.061 & 2.984 \\
IM3/residual margin (dB) & 35.6 & 44.7 \\
Tighter-repeat $|\Delta\mathrm{IIP3}|$ (dB) & 0.052 & 0.0036 \\
Longer-window $|\Delta G_V|$ (dB) & $3.43\times10^{-9}$ & $3.01\times10^{-7}$ \\
Last sideband $|\Delta\mathrm{NF}|$ (dB) & $1.31\times10^{-11}$ & $3.79\times10^{-9}$ \\
\bottomrule
\end{tabularx}
\end{table}

\reportfigure{active_compression.png}{Active mixer: conversion gain versus input power and the 1~dB compression threshold.}{fig:active-compression}

\reportfigure{active_iip3.png}{Active mixer: IF fundamentals and IM3 products versus input power per tone. Slopes are checked over $-40$ to $-30$~dBm.}{fig:active-iip3}

\reportfigure{active_twotone_spectrum.png}{Active mixer: two-tone spectrum at $-35$~dBm per tone, with 190/210~MHz fundamentals and 170/230~MHz IM3 products.}{fig:active-twotone_spectrum}

\reportfigure{passive_compression.png}{Passive mixer: conversion gain versus input power and the 1~dB compression threshold.}{fig:passive-compression}

\reportfigure{passive_iip3.png}{Passive mixer: IF fundamentals and IM3 products versus input power per tone. Slopes are checked over $-40$ to $-30$~dBm.}{fig:passive-iip3}

\reportfigure{passive_twotone_spectrum.png}{Passive mixer: two-tone spectrum at $-35$~dBm per tone, with 190/210~MHz fundamentals and 170/230~MHz IM3 products.}{fig:passive-twotone_spectrum}

\reportfigure{active_noise.png}{Active mixer: loaded SSB and DSB noise figure with 63 folded sidebands; transistor and branch-load noise are included.}{fig:active-noise}

\reportfigure{passive_noise.png}{Passive mixer: loaded SSB and DSB noise figure with 63 folded sidebands; transistor and branch-load noise are included.}{fig:passive-noise}

\subsection{Loaded Noise Figure}
The noise plots show the 63-sideband Spectre SSB and DSB results. Explicit
source-noise normalization agrees within $8.31\times10^{-5}$~dB. The last
sideband changes are negligible compared with the reported precision.
The passive circuit's lack of core DC power does not eliminate the noise
penalty of conversion loss and branch loading.

The comparison is for the stated nominal schematic models, common output
loads and ideal external sources. Layout parasitics, process variation and
LO-driver circuitry are outside the simulated circuit boundary.
